% =============================================================================
% General Introduction
% =============================================================================
\chapter*{General Introduction}
\addcontentsline{toc}{chapter}{General Introduction}
\markboth{General Introduction}{}

The digital transformation of education has accelerated considerably in recent years, driven by the widespread adoption of online learning tools and the growing expectations of both educators and students for seamless, technology-mediated collaboration.
Platforms such as Google Classroom, Moodle, and Microsoft Teams for Education have become essential components of modern academic environments.

However, despite their broad adoption, these platforms exhibit significant organizational and communicational shortcomings.
Course materials are often presented as flat, unsorted feeds that become difficult to navigate as content accumulates.
Communication channels between students and instructors are either absent or limited, and no dedicated in-platform space exists for peer-to-peer student collaboration.
Consequently, students frequently resort to external social media applications for group discussions---introducing distractions, fragmenting the learning experience, and excluding those who may not use a particular social network.

This report presents \eduspace{}, a web-based classroom management platform designed to address these gaps.
\eduspace{} structures study materials by type and chapter, integrates real-time group chat directly within each classroom, provides full-text search across all content, and supports a complete grading workflow.
The platform is architected as a microservices system.

\medskip
\noindent\textbf{Report structure.}\quad
The remainder of this report is organized as follows:

\begin{itemize}[nosep]
  \item \chapref{chap:presentation} provides the general context, objectives, development methodology, and a comparative study of existing solutions.
  \item \chapref{chap:requirements} details the requirements analysis, including actor identification, functional and non-functional requirements, and UML modeling.
  \item \chapref{chap:design} describes the system design, covering the architectural paradigm, the overall architecture, and object-oriented design artifacts.
  \item \chapref{chap:implementation} presents the implementation and deployment, including the technology stack, DevOps practices, security measures, and the user interfaces.
\end{itemize}

The report concludes with a summary of results, challenges encountered, and perspectives for future improvements.
